\section{Introducción a la Búsqueda de Ceros}

Usualmente, cuando trabajamos con ecuaciones, una de nuestras metas es encontrar el valor de cierta variable o también llamada incógnita, para esto, desarrollamos una serie de pasos o ''algoritmo'' que seguir para encontrar la respuesta. Un buen ejemplo de esto es considerar el caso de una ecuación cuadrática:

\begin{equation}
ax^2 + bx^2 +c = 0
\end{equation}

Es bien sabido que, dado los coeficientes a, b y c, existe una fórmula general que seguir, la conocida como ecuación cuadrática o \textit{Ecuación de Bashkara}:

\begin{equation}
x = \frac{-b \pm \sqrt{b^2-4ac}}{2a}
\end{equation}

Al introducir los respectivos valores de los coeficientes, podemos encontras las dos soluciones a este problema.

Sin embargo, no siempre se puede encontrar una ecuación para los valores de incógnitas en toda ecuación, un ejemplo a esto es el siguiente:

\begin{equation}
\sin x + x = 0
\end{equation}

A pesar de lo simple de nuestro problema (puesto que es la suma entre una función trigonométrica y la incógnita), ya no es posible encontrar una ecuación de forma analítica que exprese valores exactos de la soluciones para ''x'', así, la pregunta es ¿Qué podemos hacer en estos casos?

Para esto, se desarrollan los \textbf{Métodos para hallar ceros de una función}

Estos, se basan en una idea principal, convertir la ecuación en la forma:

\begin{equation}
    f(x) = 0
\end{equation}

Esto, con la finalidad de que ahora nuestro trabajo, será encontrar el valor para el cual, ''x'' cumple con ser un ''cero'' de la función.

Para encontrar este cero, es importante establecer un ''intervalo'' en donde estemos trabajando (puesto que la computadora no es capaz de analizar, en caso de trabajar en $\mathbb{R}$, la recta $]-\infty,+ \infty[$, ya que es infinita), para esto, tenemos que considerar un intervalo en donde deba existir una solución.

Sin embargo, ¿Cómo podemos asegurar una solución?

Para ello, existe el \textbf{Teorema de Bolzano}, el cual establece que:

\begin{quote}
\textit{Para una función continua en un intervalo cerrado $[a,b]$, tal que $f(a) \cdot f(b) < 0$, entonces debe existir $c \in [a,b]$ tal que $f(c) = 0$}
\end{quote}

Así, ya podemos concluir que el método con el que trabajemos esté siendo usado en un intervalo donde la solución exista.

Es importante considerar que debido a que utilizaremos métodos numéricos, la solución que encontraremos nunca será del todo correcta, sino, tan solo una aproximación a la respuesta real o analítica.

Por ello, es importante tener en cuenta \textit{¿Qué tan cerca queremos o podemos estar de nuestra respuesta real?}. Para esto, es importante tener en cuenta el uso de una medida de ''Tolerancia''.

La tolerancia, se puede definir como el intervalo más pequeño sobre el cual vamos a realizar el algoritmo de búsqueda de nuestro cero. Cuando nuestros intervalos sean suficientemente pequeños, le indicaremos al programa que detenga el algoritmo y se quede con el valor o respuesta dentro del intervalo encontrado.

Así, si la tolerancia (denotada como $\delta$) cumple que: $f(x_i) < \delta$, entonces la solución numérica encontrada será $x_i$.

Esto conlleva a la cuestión ¿Cómo podemos medir el error que se encuentra dentro de nuestro método?


% Arreglar esto!!!
Para esto, es importante analizar \textit{qué tan rápido aumenta el error}, tema del cual se verá dentro de cada método.

A continuación, exploraremos algunos de los métodos para encontrar ceros.
